\documentclass[11pt,a4paper]{article}
% ============================================================
% Shared preamble for Claude Experiments reports (pdflatex)
% Include this AFTER \documentclass in each report:
%   \documentclass[11pt,a4paper]{article}
%   % ============================================================
% Shared preamble for Claude Experiments reports (pdflatex)
% Include this AFTER \documentclass in each report:
%   \documentclass[11pt,a4paper]{article}
%   % ============================================================
% Shared preamble for Claude Experiments reports (pdflatex)
% Include this AFTER \documentclass in each report:
%   \documentclass[11pt,a4paper]{article}
%   \input{../preamble}
% ============================================================
\usepackage[T1]{fontenc}
\usepackage[utf8]{inputenc}
\usepackage[english]{babel}
\usepackage[a4paper,top=2.5cm,bottom=2.5cm,left=2.5cm,right=2.5cm]{geometry}
\usepackage{amsmath,amssymb}
\usepackage{graphicx}
\usepackage{subcaption}
\usepackage{booktabs}
\usepackage{xcolor}
\usepackage[colorlinks=true,allcolors=blue]{hyperref}
\usepackage{microtype}
\usepackage{float}
\usepackage{array}
\microtypesetup{expansion=false}

% Color scheme
\definecolor{colMicro}{HTML}{222222}
\definecolor{colExact}{HTML}{2878b5}
\definecolor{colPlacement}{HTML}{e07a1f}
\definecolor{colBoundary}{HTML}{429c74}
\definecolor{colCurrent}{HTML}{c23b55}
\definecolor{colMM}{HTML}{c23b55}

\providecommand{\SIG}{\textcolor{green!50!black}{\textbf{*}}}
\providecommand{\NS}{\textcolor{gray}{n.s.}}
\providecommand{\DRMS}{\Delta\mathrm{RMST}}
\providecommand{\ci}[2]{$[{#1},\;{#2}]$}

% ============================================================
\usepackage[T1]{fontenc}
\usepackage[utf8]{inputenc}
\usepackage[english]{babel}
\usepackage[a4paper,top=2.5cm,bottom=2.5cm,left=2.5cm,right=2.5cm]{geometry}
\usepackage{amsmath,amssymb}
\usepackage{graphicx}
\usepackage{subcaption}
\usepackage{booktabs}
\usepackage{xcolor}
\usepackage[colorlinks=true,allcolors=blue]{hyperref}
\usepackage{microtype}
\usepackage{float}
\usepackage{array}
\microtypesetup{expansion=false}

% Color scheme
\definecolor{colMicro}{HTML}{222222}
\definecolor{colExact}{HTML}{2878b5}
\definecolor{colPlacement}{HTML}{e07a1f}
\definecolor{colBoundary}{HTML}{429c74}
\definecolor{colCurrent}{HTML}{c23b55}
\definecolor{colMM}{HTML}{c23b55}

\providecommand{\SIG}{\textcolor{green!50!black}{\textbf{*}}}
\providecommand{\NS}{\textcolor{gray}{n.s.}}
\providecommand{\DRMS}{\Delta\mathrm{RMST}}
\providecommand{\ci}[2]{$[{#1},\;{#2}]$}

% ============================================================
\usepackage[T1]{fontenc}
\usepackage[utf8]{inputenc}
\usepackage[english]{babel}
\usepackage[a4paper,top=2.5cm,bottom=2.5cm,left=2.5cm,right=2.5cm]{geometry}
\usepackage{amsmath,amssymb}
\usepackage{graphicx}
\usepackage{subcaption}
\usepackage{booktabs}
\usepackage{xcolor}
\usepackage[colorlinks=true,allcolors=blue]{hyperref}
\usepackage{microtype}
\usepackage{float}
\usepackage{array}
\microtypesetup{expansion=false}

% Color scheme
\definecolor{colMicro}{HTML}{222222}
\definecolor{colExact}{HTML}{2878b5}
\definecolor{colPlacement}{HTML}{e07a1f}
\definecolor{colBoundary}{HTML}{429c74}
\definecolor{colCurrent}{HTML}{c23b55}
\definecolor{colMM}{HTML}{c23b55}

\providecommand{\SIG}{\textcolor{green!50!black}{\textbf{*}}}
\providecommand{\NS}{\textcolor{gray}{n.s.}}
\providecommand{\DRMS}{\Delta\mathrm{RMST}}
\providecommand{\ci}[2]{$[{#1},\;{#2}]$}


\title{\textbf{Metrics Reference: How to Compare MicroMacro vs.\ Full Micro}\\[0.5em]
\large Which metrics are correct, which are misleading, and why}
\author{Two-Scale CTMC Project}
\date{31 July 2026}

\begin{document}
\maketitle

\begin{abstract}
This document explains the full hierarchy of comparison metrics used in the
project, demonstrates the mathematical mechanism behind each one, and identifies
which metrics are valid under censoring and which introduce systematic bias.
The central message: \textbf{conditional mean and conditional median (completed
runs only) are biased whenever the two models have different completion rates.
RMST is the correct primary metric.}
All examples use real data from the calibrated chain-of-cliques benchmark
and the CE-01 ablation experiment.
\end{abstract}

\tableofcontents
\clearpage

%% ============================================================
\section{The Core Comparison Problem}
%% ============================================================

We compare two simulators on the same stochastic epidemic model:
\begin{itemize}
  \item \textbf{Full Micro}: exact Gillespie CTMC on the complete network.
        Ground truth.
  \item \textbf{MicroMacro (MM)}: two-scale approximation.
        Uses aggregate macro hazard $\lambda_{ij} = \beta W_{ij}\,\bar{i}_i\,\bar{s}_j$
        for inter-community transfers.
\end{itemize}

The key quantity is $T_{\mathrm{last}}$: the time when the epidemic first
reaches the last community.  In each replication the simulation runs until
$T_{\mathrm{end}}$; if no arrival has occurred by then, $T_{\mathrm{last}}$
is \emph{censored} (unknown, only known to be $> T_{\mathrm{end}}$).

The fundamental difficulty: \textbf{the two models typically have different
completion rates} (fraction of replications where $T_{\mathrm{last}} <
T_{\mathrm{end}}$), because MM spreads faster than Micro on average due to
the Jensen effect.

%% ============================================================
\section{Analyzing Method Differences: Design Before Metrics}
\label{sec:methodology}
%% ============================================================

A metric is only meaningful inside a comparison \emph{design}.  The recurring
failure mode in this project is not picking the wrong formula but asking a
metric to answer a question it cannot answer, or reading a number that silently
mixes the effect of interest with a confounder.  This section states the
principles that should govern any Micro-vs-MM comparison.  The metric catalogue
(Section~\ref{sec:catalogue}) is then just the toolbox used \emph{inside} these
principles, not a substitute for them.

\subsection{First decide which question you are asking}

Three questions look similar but demand different designs and license different
conclusions:

\begin{description}
  \item[\textbf{Q1 --- Existence.}] Does MM differ from Micro at all, beyond
        Monte-Carlo noise?  $\to$ significance of $\DRMS$ (bootstrap CI) or a
        two-sample distributional test.  Needs adequate $N$; nothing else.
  \item[\textbf{Q2 --- Magnitude and sign.}] In which direction, and by how
        much?  $\to$ $\DRMS$ and its horizon-normalised form, plus the
        completion gap.  Needs censoring-aware estimators and a fixed horizon.
  \item[\textbf{Q3 --- Mechanism / attribution.}] \emph{Why} does the difference
        arise, and which network property controls it?  $\to$ a controlled
        ablation that toggles one approximation at a time, plus a parameter
        surrogate test.  Needs confounder isolation, \emph{not} merely a large
        effect.
\end{description}

The central methodological error is answering \textbf{Q3 with a Q2 experiment}:
reporting a large, uncontrolled $\DRMS$ and \emph{narrating} a mechanism
(``premature leakage'') that the design never isolated.  A magnitude is not a
mechanism, and a sign that flips across scenarios (as in the $\beta$-ratio
sweep) is evidence that confounders, not the named mechanism, drive the number.

\subsection{The isolation principle}

To attribute a difference to a cause, every \emph{other} cause must be held
fixed.  Concretely:
\begin{itemize}
  \item \textbf{Paired seeds.}  Run Micro and every MM variant on the same
        network realisation and the same micro-layer RNG stream.  The paired
        difference $\delta_r = T_r^{MM}-T_r^{Micro}$ then removes between-run
        variance, which otherwise dominates the CI.  CE-01 does this; the
        chain-of-cliques sweep does not.
  \item \textbf{One factor at a time --- for attribution.}  To attribute error
        to a \emph{mechanism}, use the ablation ladder (Section~\ref{sec:ladder})
        that toggles one approximation per rung.  To map \emph{parameter
        dependence}, vary parameters on a grid but \emph{record every covariate}
        so its effect can be regressed out; never read a single one-factor sweep
        as causal.
  \item \textbf{Commensurable conditions.}  Two scenarios are comparable only if
        their horizon and completion regime match
        (Sections~\ref{sec:cens},~\ref{sec:horizon}).  Otherwise the numbers do
        not live on the same axis.
\end{itemize}

\subsection{Precondition 1: report completion, then use censoring-aware estimators}
\label{sec:cens}

Report $p_{\mathrm{Micro}}$, $p_{\mathrm{MM}}$ and their gap \emph{before} any
arrival-time number.  Once the gap exceeds $\approx 5$--$10\%$, every
completed-runs-only metric---conditional mean, conditional median, $W_1$,
KS---is selection-biased and must \emph{not} carry the headline
(Section~\ref{sec:bias} shows this producing the wrong sign).  Use RMST /
survival curves, which include censored runs at their known lower bound.  When
horizons cannot be matched, fall back to a Kaplan--Meier survival comparison or
re-truncate both models to a common completion level before comparing.

\subsection{Precondition 2: horizon commensurability}
\label{sec:horizon}

RMST depends on $T_{\mathrm{end}}$, so a raw $\DRMS$ is meaningful only
\emph{within} a fixed horizon.  The $\beta$-ratio sweep used
$T_{\mathrm{end}}\in[2000,20000]$; comparing its raw shifts across $r$ compares
different axes, and only the \emph{sign} and horizon-normalised quantities
survive.  Rule: within a scenario report the raw shift; \emph{across} scenarios
report $\DRMS/T_{\mathrm{end}}$, rel-$W_1$, or the completion-probability gap
$\Delta p = p_{\mathrm{MM}}-p_{\mathrm{Micro}}$.

\subsection{Choose a normalization independent of the effect}
\label{sec:norm}

\mbox{rel-}$W_1 = W_1 / \mathbb{E}[T^{Micro}]$ divides by a quantity that lives
on the \emph{same} time scale as the numerator, so it shrinks large absolute
errors ``for free'': a longer chain inflates the denominator as fast as the
numerator, and $3.5\%$ can look reassuring by construction.  Always pair it with
a scale that is independent of the mean:
\[
  \frac{W_1}{\mathrm{IQR}(T^{Micro})}
  \qquad\text{or}\qquad
  \frac{W_1}{\Delta T_{\mathrm{hop}}},\quad
  \Delta T_{\mathrm{hop}} = \frac{\mathbb{E}[T_{\mathrm{last}}]}{L-1}.
\]
A $W_1$ that is $3.5\%$ of the mean but $40\%$ of a single per-hop increment is
not small in any operationally relevant sense.

\subsection{Separate systematic signal from stochastic noise}

Between-run (phase) noise dominates unpaired comparisons and the pointwise
$I_{\mathrm{tot}}(t)$ band.  Consequences:
\begin{itemize}
  \item Prefer the \emph{paired} $\delta_i$ for Q1/Q3: it measures only the
        approximation error for a shared realisation, removing phase noise.
  \item Read the $I_{\mathrm{tot}}(t)$ IQR as ``how variable is the epidemic'',
        not ``how biased is MM'' (Section~\ref{sec:itot}).
  \item A wide band around two nearly coincident medians is \emph{not} evidence
        of a large approximation error.
\end{itemize}

\subsection{Mechanism attribution: the ablation ladder}
\label{sec:ladder}

Attribution needs a ladder whose rungs each toggle exactly one approximation, so
the increment in $\DRMS$ between adjacent rungs is charged to that one change:

\begin{center}
\begin{tabular}{llp{6.2cm}}
\toprule
Rung & Hazard $\times$ Placement & What the increment isolates \\
\midrule
Full Micro                 & exact / exact     & ground truth ($\DRMS=0$) \\
Oracle E/E                 & exact / exact-target & numerical/synchronisation only \\
Oracle E/U                 & exact / uniform   & \textbf{import placement} error \\
Aggregate A/B              & aggregate / boundary & \textbf{rate abstraction} error \\
MicroMacro A/U             & aggregate / uniform & placement $+$ rate (deployed model) \\
\bottomrule
\end{tabular}
\end{center}

CE-01 run this way attributes the bulk of the error to \emph{rate abstraction},
with placement significant only under an exact hazard.  \textbf{Caveat:} the
oracle rungs read the exact edge list and node states; they are \emph{diagnostic
upper bounds}, not deployable models---a fielded MM may know only aggregate
parameters, never the exact structure.

\subsection{From attribution to a deployable parameter surrogate}
\label{sec:surrogate}

Once the ladder localises the error to a specific approximation, the deployable
question is whether a function of \emph{known parameters}---spectral gap
$\lambda_2$, mean degree, boundary-degree statistics---can reproduce the oracle
correction \emph{without} the exact structure.  The correct test is not
``does error correlate with $\lambda_2$ across networks'' (that correlation is
confounded by everything that co-varies with $\lambda_2$).  It is: \emph{does
injecting the surrogate move MM up the ladder toward the oracle rung, with all
other factors fixed?}  To support this, a parameter sweep must record
$\lambda_2$, degree, boundary statistics and the horizon per realisation, so the
surrogate can be fitted and validated out-of-sample rather than fitted to the
same corner cases that motivated it.

\subsection{Disclose the calibration}
\label{sec:calib}

Every current bias number assumes $\beta_{\mathrm{macro}}=\beta$ and
$W_{ij}=k$ set \emph{by fiat}.  Part of $\DRMS$ is then removable by tuning
$\lambda_{ij}$ (\emph{calibration} error) and part is irreducible given only
aggregate state (\emph{structural} error).  Report the oracle-calibrated shift
as the structural floor and label every uncalibrated number as such; otherwise a
reader cannot tell a property of the method from an artefact of an arbitrary
choice.

\subsection{Include a negative control}
\label{sec:negctrl}

A credible difference analysis shows the effect \emph{vanishing} where theory
says it should.  The clique--bridge--leaf topology ($n=100$, $k=2$) is the
confounder-free control; report its $W_1$ and KS \emph{with numbers}, even when
null.  An effect that is quantified only in corner cases (small $k$, slow
spread, large $\beta_{\mathrm{out}}$) and left qualitative in the clean case is
not yet established.

\subsection{Checklist}

\begin{enumerate}
  \item State the question: Q1, Q2, or Q3.  Match the design to it.
  \item Report completion first; if the gap $>10\%$, drop completed-only metrics
        from the headline.
  \item Fix the horizon within a comparison; normalise (and disclose $T_{\mathrm{end}}$)
        across comparisons.
  \item Add an effect-independent normalisation next to rel-$W_1$.
  \item Use paired seeds and paired $\delta_i$ to strip phase noise.
  \item For mechanism claims, show the ablation ladder, not a single big shift.
  \item For parameter claims, hold covariates fixed or regress them out, and
        validate the surrogate out-of-sample.
  \item Disclose calibration; separate structural from calibration error.
  \item Include a negative control with numbers.
\end{enumerate}

\begin{table}[H]
\centering
\renewcommand{\arraystretch}{1.3}
\begin{tabular}{p{6.2cm}p{7.4cm}}
\toprule
\textbf{Pitfall} & \textbf{Correct practice} \\
\midrule
Conditional mean/median across unequal completion & Report completion; use RMST / survival (Sec.~\ref{sec:cens}) \\
Raw $\DRMS$ compared across different $T_{\mathrm{end}}$ & Normalise cross-scenario; raw only within horizon (Sec.~\ref{sec:horizon}) \\
rel-$W_1$ read as ``small'' by itself & Add $W_1/\mathrm{IQR}$ or $W_1/\Delta T_{\mathrm{hop}}$ (Sec.~\ref{sec:norm}) \\
Mechanism narrated from one large shift & Ablation ladder, one toggle per rung (Sec.~\ref{sec:ladder}) \\
Error ``explained'' by correlating with $\lambda_2$ & Fix covariates / regress out; validate surrogate (Sec.~\ref{sec:surrogate}) \\
Bias numbers with $\beta_{\mathrm{macro}}=\beta$ unlabelled & Report structural floor; mark uncalibrated (Sec.~\ref{sec:calib}) \\
Clean case left qualitative & Negative control with $W_1$/KS numbers (Sec.~\ref{sec:negctrl}) \\
\bottomrule
\end{tabular}
\caption{Mapping from common analysis pitfalls to the corresponding principle.}
\label{tab:pitfalls}
\end{table}

%% ============================================================
\section{Metric Catalogue}
\label{sec:catalogue}
%% ============================================================

\subsection{Completion Rate}

\[
  p_{\mathrm{complete}} = \frac{\#\{T_{\mathrm{last}} < T_{\mathrm{end}}\}}{N}
\]

\textbf{What it measures:} whether the epidemic traverses the entire chain
within the simulation horizon.  Affected by both the inherent stochasticity of
the epidemic and by the approximation bias of MM.

\textbf{Example (clique baseline, 100 reps):}
\[
  p_{\mathrm{Micro}} = 0.68,\quad p_{\mathrm{MM}} = 0.82.
\]
MM completes 14\% more runs---direct evidence that MM spreads faster on average.

\textbf{Correct use:} Report it first.  A large gap ($|p_{\mathrm{MM}} -
p_{\mathrm{Micro}}| > 0.1$) is already informative and means all other
completed-run metrics will be biased (see Section~\ref{sec:bias}).

\textbf{Limitation:} It is a single binary threshold; it does not quantify
\emph{how much} faster or slower.

\subsection{Conditional Mean and Median (completed runs only)}
\label{sec:bias}

\[
  \hat\mu^{\mathrm{cond}} = \frac{1}{|\mathcal{C}|}\sum_{r\in\mathcal{C}} T_r,
  \quad
  \mathcal{C} = \{r : T_r < T_{\mathrm{end}}\}.
\]

\textbf{What it measures:} the average (or median) arrival time
\emph{among runs that happened to complete}.

\paragraph{Why it is biased when completion rates differ.}
Suppose Jensen holds: $F^{MM}(t) \geq F^{Micro}(t)$ for all $t$
(MM is stochastically faster, arrives earlier).  The conditional median is
\[
  m^{Micro} = F^{-1}(0.5 \cdot p_{\mathrm{Micro}})
             = F^{-1}(0.34),
  \quad
  m^{MM}    = F^{-1}(0.5 \cdot p_{\mathrm{MM}})
             = F^{-1}(0.41).
\]
We compare $m^{MM}$ to $m^{Micro}$ but they correspond to \emph{different
percentiles} of the unconditional distribution (41st vs.\ 34th).
Even when $F^{MM}(t) \geq F^{Micro}(t)$, the 41st percentile of Micro can
easily exceed the 41st percentile of MM, making $m^{MM} > m^{Micro}$---
the opposite of what we expect.

\paragraph{Concrete example.}
Clique baseline, 100 reps:
\[
  m_{\mathrm{cond}}^{Micro} = 1949.4,\quad
  m_{\mathrm{cond}}^{MM}    = 2020.2,\quad
  \text{shift} = +70.8 \quad \text{(MM appears later!)}.
\]
But the censoring-correct RMST shift is $-75.6$ (MM is earlier).
The $+70.8$ is an artifact: MM's extra 14\% completions are slow runs near
$T_{\mathrm{end}} = 2500$ that pull the MM conditional median upward.

\textbf{Verdict: do not use conditional mean or median as the primary metric
when completion rates differ.}  They can show the wrong sign.

\subsection{Restricted Mean Survival Time (RMST)}
\label{sec:rmst}

\[
  \mathrm{RMST}
  = \mathbb{E}\!\left[\min(T_{\mathrm{last}},\, T_{\mathrm{end}})\right]
  = \int_0^{T_{\mathrm{end}}} P(T_{\mathrm{last}} > t)\, dt.
\]

\textbf{What it measures:} the expected time until arrival, capped at
$T_{\mathrm{end}}$.  Censored runs contribute $T_{\mathrm{end}}$ (their lower
bound), so all runs are included without discarding information.
Equivalently, RMST is the area under the survival curve $S(t) =
P(T_{\mathrm{last}} > t)$ from $0$ to $T_{\mathrm{end}}$.

\textbf{Why it handles censoring correctly.}
For a censored run ($T_r > T_{\mathrm{end}}$), we know the true arrival
time is at least $T_{\mathrm{end}}$.  Using $\min(T_r, T_{\mathrm{end}}) =
T_{\mathrm{end}}$ as the contribution is the correct likelihood-consistent
choice for a right-censored observation.

\textbf{Signed shift:}
\[
  \DRMS = \mathrm{RMST}^{MM} - \mathrm{RMST}^{Micro}.
\]
\begin{itemize}
  \item $\DRMS < 0$: MM arrives \emph{earlier} (premature propagation).
  \item $\DRMS > 0$: MM arrives \emph{later} (over-conservative).
  \item $\DRMS = 0$: no systematic bias at the chosen horizon.
\end{itemize}

\textbf{Connection to Jensen.}
Jensen's inequality implies $S^{Micro}(t) \geq S^{Macro}(t)$ for all $t$,
hence $\DRMS \leq 0$: integrating the survival inequality gives
$\mathrm{RMST}^{MM} \leq \mathrm{RMST}^{Micro}$.  RMST is therefore
the correct metric to \emph{test} the Jensen prediction.

\paragraph{How RMST is computed in this project (code).}
From \texttt{extrascripts/benchmark\_censoring\_summary.py}:
\begin{verbatim}
restricted_micro = np.where(np.isfinite(micro), micro, horizon)
restricted_macro = np.where(np.isfinite(macro), macro, horizon)
rmst_shift = np.mean(restricted_macro - restricted_micro)
\end{verbatim}
Censored runs ($T_r = \mathrm{NaN}$ or $\infty$) are replaced by \texttt{horizon}
($T_{\mathrm{end}}$), then the shift is the mean of the pair-wise differences
over all $N$ replications.  This is equivalent to the area difference between
survival curves estimated by the empirical CDF.

\paragraph{Bootstrap confidence interval.}
The CI is obtained by resampling the full table of $(T_r^{Micro}, T_r^{MM})$
pairs (including censored rows) with replacement, 1000 times.  Pairs are
resampled jointly (not independently) to preserve within-run correlations if
the two models share a random seed.

\paragraph{Concrete numbers (calibrated clique benchmark, 100 reps):}
\begin{center}
\begin{tabular}{lrrrl}
\toprule
Scenario & $\mathrm{RMST}^{Micro}$ & $\mathrm{RMST}^{MM}$ & $\DRMS$ & 95\% CI \\
\midrule
baseline $k=5$  & 2089.3 & 2013.6 & $-75.6$  & $[-197.3,\; +34.8]$ \\
$k=10$          & 1326.0 & 1198.1 & $-127.9$ & $[-217.8,\; \mathbf{-39.7}]$ \\
$k=20$          & 879.5  & 738.7  & $-140.8$ & $[-190.6,\; \mathbf{-89.4}]$ \\
$L=15$          & 3342.1 & 3093.7 & $-248.4$ & $[-428.7,\; \mathbf{-51.2}]$ \\
sparse baseline & 2869.4 & 1640.5 & $-1228.9$& $[-1379.4,\; \mathbf{-1089.6}]$ \\
\bottomrule
\end{tabular}
\end{center}
All scenarios with sufficient completion show $\DRMS < 0$ (bold CI = significant).
The baseline $k=5$ CI includes zero: not because Jensen is wrong, but because
100 replications are insufficient (SE $\approx 87$, shift $\approx 0.87$ SE).

\textbf{Limitation:} RMST depends on $T_{\mathrm{end}}$.  A longer horizon
makes RMST more sensitive to slowly-completing runs.  Always report the
horizon alongside RMST.

\subsection{Wasserstein-1 Distance (W\texorpdfstring{$_1$}{1})}

\[
  W_1(P, Q) = \int_{-\infty}^{\infty} |F_P(t) - F_Q(t)|\, dt
            = \int_0^1 |F_P^{-1}(u) - F_Q^{-1}(u)|\, du.
\]

\textbf{What it measures:} the average absolute difference between the quantile
functions of two distributions.  Equivalently, the minimum expected distance
needed to ``transport'' the mass of distribution $P$ onto $Q$.

\textbf{Relative W$_1$:}
\[
  \mathrm{rel}\text{-}W_1 = \frac{W_1}{\mathbb{E}[T^{Micro}]}.
\]
This normalises by the Micro mean, giving a dimensionless measure of
distributional discrepancy.

\textbf{How it is used in this project.}
$W_1$ is computed only on \textbf{completed} runs.  Therefore it is subject
to the same censoring bias as the conditional mean when completion rates
differ.  It answers: ``Among runs that \emph{both} models completed, how
different are the arrival-time distributions?''  This is valid for
\emph{distributional shape} comparison but does not capture the
``MM completes more runs'' effect.

\textbf{Concrete numbers:}
\begin{center}
\begin{tabular}{lrrl}
\toprule
Scenario & $W_1$ & rel-$W_1$ & Note \\
\midrule
clique baseline $k=5$ & 67.0 & 3.5\% & small, CI incl.\ zero \\
clique $k=10$         & 130.1 & 9.8\% & moderate \\
clique $k=20$         & 140.8 & 16.0\% & increasing with $k$ \\
sparse baseline       & ---  & 42.6\% & huge; 99\% vs 100\% complete \\
\bottomrule
\end{tabular}
\end{center}

\textbf{Limitation:} $W_1$ on completed runs is biased when completion rates
differ.  For the sparse $k=1$ scenario, Micro completes only 2\% and MM
completes 19\%; $W_1$ is not reported because the two subsets are not
comparable.

\textbf{When to prefer $W_1$ over RMST:}
\begin{itemize}
  \item When both models have 100\% completion (no censoring): $W_1$ and RMST
        give consistent information; $W_1$ is more sensitive to tail differences.
  \item When you want to compare distributional \emph{shape} (not just mean):
        $W_1$ captures quantile-by-quantile differences, not only the mean.
  \item \textbf{Not} as the primary metric when completion rates differ.
\end{itemize}

\subsection{Kolmogorov--Smirnov Statistic}

\[
  D_{\mathrm{KS}} = \sup_t |F^{Micro}(t) - F^{MM}(t)|.
\]

\textbf{What it measures:} the maximum vertical distance between the two CDFs.
The associated $p$-value tests whether the two samples could have come from
the same distribution.

\textbf{Limitation:} KS is a \emph{qualitative} test (is there a difference?)
not a \emph{quantitative} measure (how large is the difference?).  It is also
computed on completed runs only.  In this project it is used as a consistency
check, not a primary metric.

\subsection{Per-Hop Arrival-Time Metrics: W\texorpdfstring{$_1$}{1}(T\texorpdfstring{$_i$}{i}) and Paired Error}

All metrics above operate on the \emph{terminal} arrival time $T_{\mathrm{last}}$.
The per-hop analysis decomposes the error layer by layer.

\textbf{Per-hop Wasserstein $W_1(T_i)$:}
Compute the $W_1$ distance between the distribution of arrival times
$T_i^{Micro}$ and $T_i^{MM}$ for community $i = 0,\ldots,L-1$
(over all replications that reached community $i$).
This answers: ``Where in the chain does the distributional error peak?''

\textbf{Paired cumulative error:}
\[
  \delta_i = T_i^{MM} - T_i^{Micro} \quad \text{(same replication)}.
\]
This is computed only for runs where both models reached community $i$
(paired observation).  By using the same random seed, within-run correlations
are preserved and the estimator has much lower variance than between-run
comparison.

\textbf{Why the paired error is the most sensitive metric.}
Because between-run variability dominates the RMST uncertainty,
the paired comparison removes it entirely: each $\delta_i$ measures only
the \emph{approximation error} for that run, not the natural stochastic
variability.  The cost is censoring: at the last community ($i=9$) with
68\% Micro completion, only 68 pairs exist, and the missing 32 are the
\emph{slowest} Micro runs---exactly those with the largest $|\delta_i|$.

\textbf{Example (calibrated clique baseline):}
The median paired error at layer 7 is $\approx -200$ (MM arrives
200 time units earlier); at layer 9 it reverts toward 0 due to
the censoring of slow Micro runs.  The IQR is $\pm 400$---large, because
each run experiences different local dynamics.

%% ============================================================
\section{The I\texorpdfstring{$_{\mathrm{tot}}$}{tot}(t) Plot: Phase Noise vs.\ Approximation Error}
\label{sec:itot}
%% ============================================================

The total-infected count at time $t$:
\[
  I_{\mathrm{tot}}(t) = \sum_{c=0}^{L-1} I_c(t).
\]
In the chain-of-cliques setup with $n=100$ nodes per community, $I_{\mathrm{tot}}$
is approximately a step function: each community contributes a sigmoid step
of height $n$.

\paragraph{What the IQR band on $I_{\mathrm{tot}}(t)$ actually measures.}
At a fixed time $t$, the pointwise IQR across replications reflects:
\begin{enumerate}
  \item \textbf{Phase noise} ($\approx 90\%$ of width): natural variability
        in how many communities have been infected by time $t$.
        Different runs are at different ``steps'' of the staircase.
  \item \textbf{Approximation error} ($\approx 10\%$ of width): the systematic
        difference between Micro and MM in how fast communities are infected.
\end{enumerate}

Because phase noise dominates the IQR width, the $I_{\mathrm{tot}}(t)$ silhouette is
\textbf{not a useful measure of approximation quality}.  Comparing the widths of
Micro and MM IQR bands tells you about natural epidemic variability, not about
the bias of MM.

\paragraph{When the $I_{\mathrm{tot}}(t)$ plot is useful.}
\begin{itemize}
  \item \textbf{Sanity check}: if the two medians are dramatically different
        (as in the sparse baseline: $\Delta I \approx 400$ at $t=1500$),
        that is direct evidence of a large bias.
  \item \textbf{Qualitative regime identification}: the step structure shows
        how many communities are typically infected at each time.
  \item \textbf{Not for}: quantifying small-to-moderate biases (clique case),
        or comparing distributional shapes.
\end{itemize}

%% ============================================================
\section{Metric Hierarchy and Recommendations}
%% ============================================================

\begin{table}[H]
\centering
\renewcommand{\arraystretch}{1.3}
\begin{tabular}{p{3.4cm}p{2.1cm}p{4.0cm}p{4.5cm}}
\toprule
\textbf{Metric} & \textbf{Handles censoring?} & \textbf{Use when} & \textbf{Do not use when} \\
\midrule
Completion rate          & n/a        & Always report first                 & --- \\
RMST / $\DRMS$           & \textbf{Yes} & Primary metric; always valid       & --- \\
$W_1(T_{\mathrm{last}})$ & No         & Both models $\approx100\%$ complete & Completion rates differ $>5\%$ \\
Cond.\ mean shift        & No         & Completion rates $\approx$ equal    & Completion rates differ \\
Cond.\ median shift      & No         & Completion rates $\approx$ equal    & Completion rates differ \\
Paired $\delta_i$        & Partial    & Sensitive error detection per hop   & Last layer (censoring of slow runs) \\
$W_1(T_i)$ per hop       & No         & Error growth through chain          & First/last layers with few completions \\
$I_{\mathrm{tot}}(t)$    & n/a        & Regime identification, sanity check & Quantifying approximation error \\
KS statistic             & No         & Qualitative test only               & Primary quantitative comparison \\
\bottomrule
\end{tabular}
\caption{Metric hierarchy.  \textbf{Bold}: recommended primary metric.}
\label{tab:hierarchy}
\end{table}

\paragraph{Recommended reporting order for each scenario.}
\begin{enumerate}
  \item Completion rates $p_{\mathrm{Micro}}$, $p_{\mathrm{MM}}$ and their
        difference.
  \item $\DRMS$ with 95\% bootstrap CI.  Use this to test Jensen and to
        assess practical bias.
  \item rel-$W_1(T_{\mathrm{last}})$ \emph{only if} both completion rates
        $\geq 0.9$.  Provides distributional shape information beyond the mean.
  \item Per-hop paired error $\delta_i$: median and IQR across replications.
        For identifying where in the chain the error accumulates.
  \item $I_{\mathrm{tot}}(t)$ median: for visual context only, not for
        quantitative conclusions.
\end{enumerate}

%% ============================================================
\section{Worked Example: CE-01 Ablation}
%% ============================================================

CE-01 uses 100 \emph{paired} replications (same network and seed for all
variants); $T_{\mathrm{end}} = 4000$; all five variants complete 100\%
of runs (no censoring).  This is the ideal setting: all metrics are valid.

\paragraph{Primary result table.}
\begin{center}
\begin{tabular}{lrrrl}
\toprule
Variant & Completion & $\mathbb{E}[T_{\mathrm{last}}]$ & $\DRMS$ & 95\% CI \\
\midrule
Full Micro         & 1.00 & 2936.7 & $0$       & --- \\
Oracle E/E         & 1.00 & 2944.3 & $+7.5$    & $[-93.6,\;+120.3]$ \\
Oracle E/U         & 1.00 & 2775.5 & $-161.2$  & $[-266.8,\;-48.8]$ \\
Aggregate A/B      & 1.00 & 2735.5 & $-201.2$  & $[-302.4,\;-106.7]$ \\
MicroMacro (A/U)   & 1.00 & 2715.9 & $-220.8$  & $[-307.9,\;-120.2]$ \\
\bottomrule
\end{tabular}
\end{center}

Because all variants complete 100\% of runs and share the same seeds,
$\DRMS = \mathbb{E}[T^{\mathrm{variant}}] - \mathbb{E}[T^{Micro}]$ exactly
(no censoring, no between-run confounding).  The conditional mean and RMST
are identical here.

\paragraph{Why RMST still matters even at 100\% completion.}
The RMST formulation generalises cleanly to the censored case.
Using RMST throughout---even when it reduces to the plain mean---ensures
the same code and the same interpretation across all experiments.
This avoids introducing separate code paths that could silently compute
the wrong thing when some runs are censored.

%% ============================================================
\section{Statistical Power and Sample Size}
%% ============================================================

A common question: how many replications are needed to detect a given
$\DRMS$ with 80\% power at the 5\% significance level?

For a one-sample $t$-test on $\Delta_r = \min(T_r^{MM}, T_{\mathrm{end}}) -
\min(T_r^{Micro}, T_{\mathrm{end}})$ with standard deviation $\sigma_\Delta$:
\[
  N \approx \left(\frac{(z_{0.025} + z_{0.20})\, \sigma_\Delta}{\delta}\right)^2
           = \left(\frac{2.80\, \sigma_\Delta}{\delta}\right)^2,
\]
where $\delta = |\DRMS|$ is the target effect size.

\paragraph{Calibrated clique baseline ($k=5$).}
From bootstrap: $\sigma_\Delta \approx 590$ (std of pairwise differences),
$|\DRMS| = 75.6$.  Required $N \approx (2.80 \times 590 / 75.6)^2 \approx 480$.
Current 100 replications give power $\approx 25\%$---consistent with the CI
including zero.

\paragraph{Calibrated clique $k=10$.}
$\sigma_\Delta \approx 450$, $|\DRMS| = 127.9$.
Required $N \approx (2.80 \times 450 / 127.9)^2 \approx 97$.
100 replications are just sufficient; CI excludes zero.

\paragraph{CE-01 (paired).}
Pairing the replications (same seed) reduces $\sigma_\Delta$ dramatically
(between-run variance is removed).  100 paired replications are sufficient
for all four contrasts with $|\DRMS| \geq 160$.

\textbf{Practical rule of thumb:} for the chain benchmark without pairing,
use at least $N = 300$ replications when $|\DRMS| < 100$.

%% ============================================================
\section{Summary}
%% ============================================================

\begin{enumerate}
  \item \textbf{RMST and $\DRMS$} are the correct primary metrics.
        They handle censoring, connect directly to the Jensen inequality,
        and have well-understood statistical properties.

  \item \textbf{Conditional mean/median} (completed runs only) introduces
        selection bias whenever completion rates differ.  It can show the
        \emph{wrong sign} (e.g., clique baseline: conditional median $= +70.8$
        while RMST $= -75.6$).

  \item \textbf{W$_1$ distance} is valid for distributional shape comparison
        when both models have near-100\% completion.  Use it alongside RMST,
        not instead of it.

  \item \textbf{$I_{\mathrm{tot}}(t)$ pointwise plots} reflect phase noise,
        not approximation error.  Useful for visual context (especially
        sparse case) but not for quantifying bias.

  \item \textbf{Per-hop paired error $\delta_i$} is the most sensitive
        diagnostic: it removes between-run variance by using paired seeds.
        Interpret the last-layer value cautiously due to censoring of
        slow Micro runs.

  \item \textbf{100 replications} are insufficient for the clique baseline
        ($k=5$) to detect the Jensen bias at 5\% significance.
        Need $\approx 400$--$500$ replications or paired seeds.
\end{enumerate}

\end{document}
